% =========================================================
%  General Formularium – Mathematics
%  Chapter 2: Geometry & Trigonometry
%  Converted from "General Formularium (ENG) - V2.docx"
%  Reviewed and corrected on 21 Oct 2026.
% =========================================================
\documentclass[11pt]{article}
\usepackage[a4paper,margin=2cm]{geometry}
\usepackage{amsmath,amssymb}
\usepackage[dvipsnames]{xcolor}
\setlength{\parindent}{0pt}
\newcommand{\sect}[2]{\section*{\color{#1}#2}}
\newcommand{\sub}[1]{\par\medskip\textbf{#1}\par}
\newcommand{\expl}[1]{{\color{gray}\itshape #1}\par}
\newcommand{\warn}[1]{{\color{BrickRed}\textbf{!} #1}}
\definecolor{cGeo}{HTML}{2E8B57}

\begin{document}

\sect{cGeo}{Geometry \& Trigonometry}
\expl{Areas and volumes of common shapes, and the trigonometric identities.}

% ---------- Areas of plane figures ----------
\sub{Area of a rectangle}
\[ A = b \cdot h \]

\sub{Area of a parallelogram}
\[ A = b \cdot h \]

\sub{Area of a triangle}
\[ A = \tfrac{1}{2}\, b \cdot h \]

\sub{Area of a trapezoid}
\[ A = \frac{b + B}{2}\cdot h \]
\expl{$b$ and $B$ are the two parallel sides.}

\sub{Area of a circle}
\[ A = \pi r^2 \]

\sub{Area of an ellipse}
\[ A = \pi a b \]
\expl{$a$ and $b$ are the semi-axes.}

% ---------- Surface areas of solids ----------
\sub{Surface area of a cube}
\[ A = 6s^2 \]

\sub{Surface area of a rectangular prism (box)}
\[ A = 2(lb + lh + bh) \]

\sub{Surface area of a triangular prism}
\[ A = 2A_{\text{triangle}} + p_{\text{triangle}}\cdot l \]
\expl{$p_{\text{triangle}}$ = perimeter of the triangle, $l$ = length of the prism.}

\sub{Surface area of a pyramid}
\[ A = A_{\text{base}} + \textstyle\sum A_{\text{triangles}} \]

\sub{Surface area of a cylinder}
\[ A = 2\pi r h + 2\pi r^2 \]

\sub{Surface area of a cone}
\[ A = \pi r s + \pi r^2 \]
\expl{$s$ = slant height (schuine zijde).}

\sub{Surface area of a sphere}
\[ A = 4\pi r^2 \]

% ---------- Volumes ----------
\sub{Volume of a cube}
\[ V = s^3 \]

\sub{Volume of a rectangular prism (box)}
\[ V = l \cdot b \cdot h \]

\sub{Volume of a cylinder}
\[ V = \pi r^2 h \]

\sub{Volume of a pyramid}
\[ V = \tfrac{1}{3} A_{\text{base}}\, h \]

\sub{Volume of a cone}
\[ V = \tfrac{1}{3}\pi r^2 h \]

\sub{Volume of a sphere}
\[ V = \tfrac{4}{3}\pi r^3 \]

% ---------- Trigonometry ----------
\sub{Pythagorean theorem}
\[ a^2 + b^2 = c^2 \]
\expl{$c$ is the hypotenuse.}

\sub{Trigonometric ratios in a right triangle (SOH CAH TOA)}
\[ \sin\alpha = \frac{\text{opposite}}{\text{hypotenuse}} \qquad \cos\alpha = \frac{\text{adjacent}}{\text{hypotenuse}} \qquad \tan\alpha = \frac{\text{opposite}}{\text{adjacent}} \]

\sub{Fundamental identity}
\[ \sin^2\alpha + \cos^2\alpha = 1 \]

\sub{Trigonometric ratios}
\[ \tan\alpha = \frac{\sin\alpha}{\cos\alpha} \qquad \cot\alpha = \frac{\cos\alpha}{\sin\alpha} \qquad \sec\alpha = \frac{1}{\cos\alpha} \qquad \csc\alpha = \frac{1}{\sin\alpha} \]
\expl{Valid where the denominator is not 0.}

\sub{Cosine rule}
\[ a^2 = b^2 + c^2 - 2bc\cos\alpha \]

\sub{Sine rule}
\[ \frac{a}{\sin\alpha} = \frac{b}{\sin\beta} = \frac{c}{\sin\gamma} \]

\sub{Sum and difference formulas}
\[ \sin(\alpha \pm \beta) = \sin\alpha\cos\beta \pm \cos\alpha\sin\beta \]
\[ \cos(\alpha \pm \beta) = \cos\alpha\cos\beta \mp \sin\alpha\sin\beta \]
\[ \tan(\alpha \pm \beta) = \frac{\tan\alpha \pm \tan\beta}{1 \mp \tan\alpha\tan\beta} \]

\sub{Double-angle formulas}
\[ \sin 2\alpha = 2\sin\alpha\cos\alpha \qquad \tan 2\alpha = \frac{2\tan\alpha}{1 - \tan^2\alpha} \]
\[ \cos 2\alpha = \cos^2\alpha - \sin^2\alpha = 2\cos^2\alpha - 1 = 1 - 2\sin^2\alpha \]

\sub{Power-reduction formulas (Carnot)}
\[ \sin^2\alpha = \frac{1 - \cos 2\alpha}{2} \qquad \cos^2\alpha = \frac{1 + \cos 2\alpha}{2} \]

\sub{t-formulas (half-angle substitution)}
\[ t = \tan\frac{\alpha}{2}: \qquad \sin\alpha = \frac{2t}{1 + t^2} \qquad \cos\alpha = \frac{1 - t^2}{1 + t^2} \qquad \tan\alpha = \frac{2t}{1 - t^2} \]

\sub{Product-to-sum formulas (inverse Simpson)}
\[ \cos\alpha\cos\beta = \tfrac{1}{2}\left[\cos(\alpha - \beta) + \cos(\alpha + \beta)\right] \]
\[ \sin\alpha\sin\beta = \tfrac{1}{2}\left[\cos(\alpha - \beta) - \cos(\alpha + \beta)\right] \]
\[ \sin\alpha\cos\beta = \tfrac{1}{2}\left[\sin(\alpha + \beta) + \sin(\alpha - \beta)\right] \]

\sub{Sum-to-product formulas (Simpson)}
\[ \sin p + \sin q = 2\sin\frac{p+q}{2}\cos\frac{p-q}{2} \qquad \sin p - \sin q = 2\cos\frac{p+q}{2}\sin\frac{p-q}{2} \]
\[ \cos p + \cos q = 2\cos\frac{p+q}{2}\cos\frac{p-q}{2} \qquad \cos p - \cos q = -2\sin\frac{p+q}{2}\sin\frac{p-q}{2} \]

\sub{Solving basic trigonometric equations}
\[ \sin x = \sin\alpha \iff x = \alpha + k2\pi \;\text{ or }\; x = \pi - \alpha + k2\pi \]
\[ \cos x = \cos\alpha \iff x = \alpha + k2\pi \;\text{ or }\; x = -\alpha + k2\pi \]
\[ \tan x = \tan\alpha \iff x = \alpha + k\pi \qquad \cot x = \cot\alpha \iff x = \alpha + k\pi \]
\expl{$k \in \mathbb{Z}$.}

\sub{Equal angles (periodicity)}
\[ \sin(\alpha + 2k\pi) = \sin\alpha \qquad \cos(\alpha + 2k\pi) = \cos\alpha \]

\sub{Complementary angles}
\[ \sin\left(\frac{\pi}{2} \pm \alpha\right) = \cos\alpha \qquad \cos\left(\frac{\pi}{2} \pm \alpha\right) = \mp\sin\alpha \]

\end{document}
